\documentclass{beamer}
\usepackage[utf8]{inputenc}
\usetheme{default}
\title{Active Networks}
\author{Arseny Kurnikov\\
Hasan Islam\\
Harri Hämäläinen}
\begin{document}


\begin{frame}
\titlepage
\end{frame}

\begin{frame}[fragile]{Goals of Active Networks}
	\begin{itemize}
		\item Create a dynamic and extensible communications environment that allow application specific customization.
		\item Introduce Programmable open nodes with the ability to deploy programs dynamically into node engines.
        \item Some use cases:
        \begin{itemize}
            \item Data distribution and aggregation
            \item Active network management and measurement
            \item Network security and protection
        \end{itemize}
	\end{itemize}
	
\end{frame}

\begin{frame}[fragile]{Active Networks}
    \begin{itemize}
        \item Applications can add code into packets which can be executed by
        the communication hardware to provide dynamic behaviour tailored for the
        application needs.
		\item The Discrete Approach: programs are injected into active network node separately from actual data packet
		\item The Integrated Approach: program is integrated into every packet of data
    \end{itemize}
    \begin{verbatim}
 +--------+------+         +--------+------+------+
 | Header | Data |   -->   | Header | Code | Data |
 +--------+------+         +--------+------+------+
    \end{verbatim}
\end{frame}

\begin{frame}{Code}
  Goals
	\begin{itemize}
		\item Mobility
    \item Safety
    \item Efficiency
	\end{itemize}
  Possible solutions
  \begin{itemize}
    \item Source code
    \item Intermediate code (Java)
    \item Binary code (platform-dependent)
    \item On-the-fly compilation
  \end{itemize}
	
\end{frame}

\begin{frame}{End-to-end argument and active networking}
    \begin{itemize}
        \item \textit{''application-specific functions ought to reside in the end hosts
        of a network rather than in intermediary nodes – provided they can be
        implemented \textbf{completely and correctly} in the end hosts''} - End-to-end
        principle
        \item On one hand the programmability of the network breaks the
        end-to-end principle
        \item On the otherhand the tailored service exactly for the application
        needs is in a sense the ''correct'' method.
    \end{itemize}
\end{frame}

\begin{frame}{Tradeoffs}
    \begin{itemize}
        \item Security
        		\begin{itemize}
        			\item who can execute code where?
        			\item how do you limit what actions can be taken?
        			\item who is allowed to send code to the routers?
        		\end{itemize}
        \item Performance
        			\begin{itemize}
        				\item monopolizing the processor by a malicious piece of code
        				\item expensive verification, while loading code into memory probably, enough to slow down the system and cause momentary performance loss
        			\end{itemize}
    \end{itemize}
\end{frame}

\begin{frame}{Sources}
    \begin{itemize}
        \item A Survey of Active Network Research, Tennenhouse et al. 1997
        \item SANDS: Specialized Active Networking for Distributed Simulation,
        Zabele et al., 2002
        \item Active Networking and End-To-End Arguments, Reed, Saltzer, and
        Clark
    \end{itemize}
\end{frame}

\end{document}
